\section{Introduction}

While matrix-level distillation faithfully captures attention patterns at training lengths, token-level alignment produces representations 5$\times$ more stable across sequence lengths---a mechanistic advantage suggesting superior generalization potential. This finding reframes Transformer-to-SSM distillation as an objective selection problem rather than purely an architectural one.

\subsection{The Long-Context Challenge}

Transformers achieve remarkable performance on natural language tasks but scale quadratically with sequence length, making long-context processing prohibitively expensive. A 32K-token sequence requires 256$\times$ more compute than a 2K sequence for attention alone. This scaling barrier has motivated extensive research into linear-time alternatives, particularly State Space Models (SSMs) like Mamba \citep{gu2023mamba}.

Rather than training SSMs from scratch---requiring billions of tokens and substantial compute---distillation offers a compelling alternative. MOHAWK \citep{bick2024mohawk} demonstrates that Phi-Mamba can match 85\% of Phi-1.5's performance using only 3B tokens ($<$1\% of typical pretraining cost). CAB \citep{wang2025cab} shows that token-level alignment avoids O($L^2$) attention map materialization entirely.

\subsection{The Overlooked Question}

Yet a fundamental question remains unexplored: \textit{which distillation objective generalizes better when target sequences exceed training lengths?}

MOHAWK employs matrix-level supervision, minimizing $\|$TeacherMixer - StudentMixer$\|$ to directly transfer attention patterns. CAB uses token-level supervision, aligning query/key projections (Q,K) to SSM parameters (B,C) via learned bridges. Both achieve competitive results at training lengths, but neither has been evaluated for length extrapolation under controlled conditions.

This gap matters because practitioners increasingly need models that handle sequences longer than their training distribution. If one objective produces representations that transfer better to unseen lengths, this provides a principled criterion for distillation design.

\subsection{Key Insight: Position-Agnostic Supervision Transfers Better}

We hypothesize that token-level objectives produce more stable representations across lengths because:

\begin{enumerate}
    \item \textbf{Token-level supervision is position-agnostic}: Aligning individual Q/K projections to B/C parameters creates supervision independent of sequence-specific relationships.
    \item \textbf{Matrix-level supervision encodes positional patterns}: Attention maps capture position-position relationships that fundamentally change with sequence length.
\end{enumerate}

Our experiments confirm this mechanism. Measuring hidden state drift (L2 distance between teacher and student representations) across sequence lengths 512--2048, we find CAB produces drift slope 0.00090506 versus MOHAWK's 0.00452495---a 5$\times$ difference. CAB's drift ratio remains bounded at 1.34; MOHAWK's is 2.02.

\subsection{Contributions}

This work makes three contributions:

First, we construct a unified Phi-Mamba framework implementing both MOHAWK and CAB objectives in a single codebase (Section~\ref{sec:methodology}). This enables the first controlled comparison of distillation objective types, eliminating confounds from implementation differences.

Second, we demonstrate that token-level distillation produces 5$\times$ more stable representations across sequence lengths (Section~\ref{sec:results}). This quantified stability gap provides mechanistic evidence that objective type fundamentally affects length generalization.

Third, we establish scope conditions for Transformer-to-SSM distillation: Phi-1.5's hard 2048-token context limit (Section~\ref{sec:experiments}) reveals that length extrapolation research requires teachers with explicit position extrapolation capabilities (RoPE, ALiBi). Our findings suggest token-level objectives may be preferred when target lengths exceed the teacher's training distribution.
